% !TeX root = ../main.tex
% 【盲审占位】非第一作者论文是否纳入本节，视所在学院研究生院细则而定。
% 正式提交时按学院要求补充完整的论文列表、作者排序、刊物/会议全称、年/卷/期/页码或 DOI。
\begin{achievelist}
    \item Cai JW, Wei YC, Wu ZT, \textbf{Peng S}, Ma KS. Inter-layer scheduling space definition and exploration for tiled accelerators[C]//ACM/IEEE International Symposium on Computer Architecture (ISCA 2023). 2023: 1--17. DOI: 10.1145/3579371.3589048.
    \item Cai JW, Wu ZT, \textbf{Peng S}, et al. Gemini: Mapping and architecture co-exploration for large-scale DNN chiplet accelerators[C]//IEEE International Symposium on High-Performance Computer Architecture (HPCA 2024). 2024: 156--171. DOI: 10.1109/HPCA57654.2024.00022.
    {\sloppy\item Cai JW, Wang X, Gao MY, \textbf{Peng S}, et al. SoMa: Identifying, exploring, and understanding the DRAM communication scheduling space for DNN accelerators[C]//IEEE International Symposium on High-Performance Computer Architecture (HPCA 2025). 2025: 533--548. DOI: 10.1109/\allowbreak HPCA61900.\allowbreak 2025.\allowbreak 00048.\par}
    \item Wei YC, Cai JW, Gao MY, \textbf{Peng S}, et al. Buffer prospector: Discovering and exploiting untapped buffer resources in many-core DNN accelerators[C]//ACM/IEEE Design Automation Conference (DAC 2025). 2025: 1--7. DOI: 10.1109/DAC63849.2025.11132440.
\end{achievelist}

% ======================================
%       正式论文请将以下内容注释或删除
% ======================================
% \vspace{\baselineskip}
% {\color{red} 用于双盲评审的论文，只列出已发表的学术论文的篇名、发表刊物名称，必须隐去各类论文检索号、期号、卷号、页码；专利号；日期等。}